\section{Pragmatics}
\label{sec:pragmatics}
\slidesrc{Pragmatics: 1--66}

We reach the third level of linguistic analysis (Section~\ref{sec:three-levels}):
after \emph{what is grammatical?} (syntax) and \emph{what does it mean?} (semantics),
\textbf{pragmatics} asks \emph{what does it do?} What does the speaker want to achieve by
saying it, and how does the context change what is understood? The first slide sets the
tone with a comic in which a character believes he has ``decoded'' what \emph{yes},
\emph{maybe} and \emph{no} really mean: the joke works because what is \emph{said} and
what is \emph{meant} can differ, which is exactly the subject of pragmatics.

This section covers:
\begin{enumerate}[nosep]
  \item what pragmatics is (Section~\ref{sec:what-is-pragmatics});
  \item communication as action: Wittgenstein, Austin, Searle, speech acts and the
        stages of communication (Section~\ref{sec:communication-as-action});
  \item whether language shapes thought: Newspeak, the Sapir--Whorf hypothesis and the
        experiments presented by Lera Boroditsky (Section~\ref{sec:language-and-thought});
  \item the language of numbers and the power of names
        (Section~\ref{sec:numbers-and-names}).
\end{enumerate}

\subsection{What Is Pragmatics}
\label{sec:what-is-pragmatics}
\slidesrc{Pragmatics: 1--3}

\begin{nbdef}{Pragmatics}
\textbf{Pragmatics} is a subfield of linguistics and semiotics that studies the ways in
which \textbf{context contributes to meaning}. It encompasses speech act theory,
conversational implicature, talk in interaction and other approaches to language
behaviour in philosophy, sociology, linguistics and anthropology.

Pragmatics studies how the transmission of meaning depends not only on the structural
and linguistic knowledge (grammar, lexicon, \ldots) of speaker and listener, but also on
\begin{itemize}[nosep]
  \item the \textbf{context of the utterance};
  \item any \textbf{pre-existing knowledge} about those involved;
  \item the \textbf{inferred intent} of the speaker;
  \item and other factors.
\end{itemize}
\end{nbdef}

\textbf{Semiotics} is the study of meaning-making: the study of sign processes
(\emph{semiosis}) and of meaningful communication. Words are one kind of sign; gestures,
pictures, traffic lights are others.

\begin{nbexample}{same words, different meanings}
``Can you pass the salt?'' Syntactically it is a question; semantically it asks about an
ability (are you able to pass the salt?). Pragmatically, at a dinner table, it is a
\emph{request}: answering ``yes'' without moving would be absurd. ``It's cold in
here'' can be a statement, a complaint, or a request to close the window, depending on
who says it to whom, and where. Only the context tells.
\end{nbexample}

\begin{nbwhy}
A chatbot that understands every word but not the intention behind it answers ``Yes, I
can'' to ``Can you tell me the time?''. Recognising intentions (\emph{intent detection})
is a central task of dialogue systems, and the reason why the agent communication
languages of Section~\ref{sec:chatbdi} label every message with the act it performs.
\end{nbwhy}

\subsection{Communication as Action}
\label{sec:communication-as-action}
\slidesrc{Pragmatics: 4--10}

These slides come from another course (on agents) and give a very quick introduction to
``communication as an action''.

\subsubsection{From true/false sentences to actions}
\slidesrc{Pragmatics: 5--8}

\begin{nbdef}{Communication}
\textbf{Communication}, in general, is the \emph{intentional} exchange of information
brought about by the production and perception of \emph{signs} drawn from a
\emph{shared system of conventional signs}.
\end{nbdef}

Two views of language:
\begin{itemize}[nosep]
  \item the \textbf{``classical'' view} (before 1953): language consists of sentences
        that are \textbf{true or false} (as in logic; this is the view behind
        model-theoretic semantics, Section~\ref{sec:semantic-parsing});
  \item the \textbf{``modern'' view} (after 1953): \textbf{language is a form of
        action}.
\end{itemize}
The turning point is marked by three philosophers.

\paragraph{Wittgenstein (1953), \emph{Philosophical Investigations}.}
According to Wittgenstein:
\begin{itemize}[nosep]
  \item the structure of reality determines the structure of language;
  \item \textbf{meaning is determined by use}, in the context of a
        \textbf{``language-game''};
  \item words are not defined by reference to the objects or things they designate in
        the external world, nor by the thoughts, ideas or mental representations one
        might associate with them, but by \textbf{how they are used} in effective,
        ordinary communication.
\end{itemize}

\begin{nbintuition}[language-games]
Think of the word \emph{``Slab!''} shouted by a builder to an assistant: it does not
describe a slab, it means ``bring me a slab'', and it means that only inside the
``game'' of building together. Every use of language is a move in some game with its own
rules: giving orders, asking, telling a joke, praying. This is the same Wittgenstein of
``meaning is use'' at the root of distributional semantics
(Section~\ref{sec:distributional-hypothesis}).
\end{nbintuition}

\begin{nbwarning}[two Wittgensteins]
The first bullet (``the structure of reality determines the structure of language'') is
the idea of the \emph{early} Wittgenstein (\emph{Tractatus}, 1921), in which a sentence is
a picture of a fact. The \emph{Philosophical Investigations} (1953) largely abandon it in
favour of meaning as use. The slide lists both; remember that the 1953 work is the one of
language-games.
\end{nbwarning}

\paragraph{Austin (1962), \emph{How to Do Things with Words}.}
\textbf{In saying something, we do something}: when a minister joins two people in
marriage saying ``I now pronounce you husband and wife'', the sentence does not describe
a marriage, it \emph{makes} it (Section~\ref{sec:course-map}). Such sentences are
\textbf{performatives}.

\paragraph{Searle (1969), \emph{Speech Act Theory}.}
\begin{itemize}[nosep]
  \item Each \textbf{illocution} (the act performed \emph{in} saying something) can be
        described in terms of what it is attempting to do. An \emph{assertion} counts as
        an undertaking that something really is the case; a \emph{question} counts as an
        attempt to elicit some information; \emph{thanking} someone counts as an
        expression of gratitude.
  \item What matters is the assumed \textbf{intent of the speaker}, the
        \textbf{intentionality} of the sentence.
\end{itemize}

\notinslides
\begin{nbdef}{Speech act: three levels (Austin)}
\begin{itemize}[nosep]
  \item \textbf{Locution}: the act of saying something, the words with their literal
        meaning (``It's cold in here'').
  \item \textbf{Illocution}: the act performed \emph{in} saying it, the speaker's
        intention (a request to close the window).
  \item \textbf{Perlocution}: the effect produced \emph{by} saying it on the hearer
        (the hearer gets up and closes the window).
\end{itemize}
\end{nbdef}

\subsubsection{Speech acts and the stages of communication}
\slidesrc{Pragmatics: 9--10}

\begin{center}
\begin{tikzpicture}[b/.style={font=\small}]
  \node[draw=nbred, rounded corners=6pt, minimum width=6.2cm, minimum height=1.1cm] (sit) {};
  \node[font=\tiny, anchor=north west, text=nbred] at (sit.north west) {SITUATION};
  \node[b, text=nbblue] (s) at (-2.0,-0.1) {Speaker};
  \node[b, text=nbred] (u) at (0.1,-0.1) {Utterance};
  \node[b, text=nbblue] (h) at (2.2,-0.1) {Hearer};
  \draw[-{Stealth}] (s) -- (u); \draw[-{Stealth}] (u) -- (h);
\end{tikzpicture}
\captionof{figure}{A speech act: a speaker produces an utterance for a hearer, inside a
situation.}
\end{center}

\textbf{Speech acts achieve the speaker's goals}. The examples come from the Wumpus world,
the cave game used for the grammar of Section~\ref{sec:grammars-parsing}:
\begin{center}
\small
\begin{tabular}{@{}l l@{}}
\toprule
\textbf{act} & \textbf{example} \\
\midrule
Inform & ``There's a pit in front of you'' \\
Query & ``Can you see the gold?'' \\
Command & ``Pick it up'' \\
Promise & ``I'll share the gold with you'' \\
Acknowledge & ``OK'' \\
\bottomrule
\end{tabular}
\end{center}
\textbf{Speech act planning} (deciding what to say to reach a goal) requires knowledge
of the \emph{situation}, of the \emph{semantic and syntactic conventions}, and of the
\emph{hearer's goals, knowledge base and rationality}.

\begin{center}
\small
\begin{tabular}{@{}l l l@{}}
\toprule
 & \textbf{stage} & \textbf{what happens} \\
\midrule
\multirow{3}{*}{\textcolor{nbblue}{speaker S}}
  & Intention & S wants to inform H that $P$ \\
  & Generation & S selects words $W$ to express $P$ in context $C$ \\
  & Synthesis & S utters words $W$ \\
\midrule
\multirow{4}{*}{\textcolor{nbblue}{hearer H}}
  & Perception & H perceives $W'$ in context $C'$ \\
  & Analysis & H infers possible meanings $P_1, \ldots, P_n$ \\
  & Disambiguation & H infers the intended meaning $P_i$ \\
  & Incorporation & H incorporates $P_i$ into its knowledge base \\
\bottomrule
\end{tabular}
\captionof{table}{The stages in communication.}
\label{tab:communication-stages}
\end{center}

\begin{nbintuition}[where things can go wrong]
Notice the primes: the hearer perceives $W'$ in context $C'$, not necessarily the same
$W$ and $C$ of the speaker (noise, mishearing, a different idea of the situation). The
analysis gives \emph{several} possible meanings, and disambiguation picks one, maybe the
wrong one. Every stage maps onto a part of NLP: generation and synthesis are natural
language generation and speech synthesis; perception is speech recognition or reading;
analysis is syntax and semantics (parsing, Sections~\ref{sec:sentence-level-syntax}
to~\ref{sec:lexical-semantics}); disambiguation is WSD and intent recognition;
incorporation is updating a knowledge base, e.g.\ an ontology
(Section~\ref{sec:ontologies}).
\end{nbintuition}

\subsection{Does Language Shape Thought?}
\label{sec:language-and-thought}
\slidesrc{Pragmatics: 11--49}

``Thoughts shape language\ldots\ but\ldots\ does language shape thought?'' If pragmatics
studies how context and intentions shape what we mean, the reverse question is natural:
does the language we speak shape how we think?

\subsubsection{Newspeak and the Sapir--Whorf hypothesis}
\slidesrc{Pragmatics: 13--16}

In George Orwell's novel \emph{Nineteen Eighty-Four} (1949), the regime designs a new
language, called \textbf{Newspeak}, to control thought:

\begin{quote}\small
``Don't you see that the whole aim of Newspeak is to narrow the range of thought? In the
end we shall make thought\-crime literally impossible, because there will be no words in
which to express it. Every concept that can ever be needed will be expressed by exactly
one word, with its meaning rigidly defined and all its subsidiary meanings rubbed out and
forgotten. [\ldots] Every year fewer and fewer words, and the range of consciousness
always a little smaller. [\ldots] The Revolution will be complete when the language is
perfect.''
\end{quote}

\begin{nbdef}{Linguistic relativity (Sapir--Whorf hypothesis)}
The hypothesis of \textbf{linguistic relativity} (also \textbf{Sapir--Whorf
hypothesis}, or Whorfianism, after Edward Sapir, 1884--1939, and Benjamin Lee Whorf,
1897--1941) holds that the structure of a language affects its speakers' world view or
cognition. It has two versions:
\begin{itemize}[nosep]
  \item the \textbf{strong} version, \textbf{linguistic determinism}: language
        \emph{determines} thought, linguistic categories limit and determine cognitive
        categories; language entirely determines the range of cognitive processes;
  \item the \textbf{weak} version, \textbf{linguistic influence}: linguistic categories
        and usage only \emph{influence} thought and decisions; language provides
        constraints in some areas of cognition, but is by no means determinative.
\end{itemize}
\end{nbdef}

\begin{nbfact}{what we know today}
Linguistic \textbf{determinism} (the strong version) is now generally agreed to be
\textbf{false}: we can think things for which our language has no word (otherwise new
words could never be invented). Research on the \textbf{weaker} forms has produced
\textbf{positive empirical evidence}: language does influence some areas of cognition.
Newspeak is the strong version, and that is why it is fiction.
\end{nbfact}

\subsubsection{How languages and cultures shape the ways we think}
\slidesrc{Pragmatics: 17--30}

The rest of this part follows a talk by \textbf{Lera Boroditsky} (Associate Professor of
Cognitive Science at UC San Diego, previously at MIT and Stanford, who studies the
relationships between mind, world and language), \emph{How languages and cultures shape
the ways we think}.

\paragraph{Languages force us to say different things.}
Take the English sentence \emph{``Trump read Ta-Nehisi's latest book''}. Depending on the
language, to say it you \emph{must} decide:
\begin{itemize}[nosep]
  \item did this happen in the past? (tense: Mandarin does not require it, English does);
  \item Donald or Melania? (in Russian the past-tense verb changes with the gender of the
        subject);
  \item was the action completed? (aspect: Russian forces a choice);
  \item did you witness this with your own eyes? (evidentiality: in Turkish, for
        example, the verb marks whether you saw it or heard about it).
\end{itemize}
Similarly, \emph{``the 7th blue house to the left of the big tree''}: seventh? blue, light
or dark (Russian has no single word for ``blue'')? left of, or north-northwest of, or
``volcano-ward'' of (some languages use absolute directions)? and which tree: banyan, red
cypress, eucalyptus? Each language makes its speakers pay attention to different
aspects of the world, all the time.

\paragraph{How many languages in the world?}
According to Ethnologue: 7\,168 languages in use (October 2023), 7\,164 (October 2024),
7\,159 (November 2025). The number is constantly in flux, because we keep learning about
languages and because languages themselves change. Roughly \textbf{44\%} of all languages
are \textbf{endangered}, often with fewer than 1\,000 users left, while the 20 largest
languages are the native tongue of more than 3.6--3.7 billion people: about 0.3\% of the
languages account for nearly half of the world's population. Another source is the
\emph{World Atlas of Language Structures} (WALS), which describes the structural
features of a sample of languages.

\paragraph{Navajo verbs of eating.}
Navajo has no single verb ``to eat'': the verb stem changes with the \emph{kind} of food
and the \emph{manner} of eating.

\begin{center}
\small
\begin{tabular}{@{}r l l l@{}}
\toprule
 & \textbf{stem} & \textbf{manner / food of consumption} & \textbf{example edibles} \\
\midrule
i.    & -y\'{\i}\'{\i}h & consume/chew \textsc{unspecified} & anything \\
ii.   & -'aa\l   & consume/chew \textsc{hard, compact} & corn, peyote, candy \\
iii.  & -chozh   & consume/chew \textsc{leafy} & lettuce, hay \\
iv.   & -gha\l   & consume/chew \textsc{meat} & meat \\
v.    & -k\'e\'ed & consume/chew \textsc{round} & bun, melon \\
vi.   & -ts'\'e\'eh & consume/chew \textsc{mushy matter} & mush, jello \\
vii.  & -dee\l   & consume/chew \textsc{plural object} & berries, eggs \\
viii. & -wol     & consume/chew \textsc{marrow} & marrow \\
ix.   & -joo\l   & consume \textsc{non-compact matter} & cotton candy \\
x.    & -t\l eeh & consume \textsc{mushy matter} & ice cream \\
xi.   & -m\'a\'a' & consume/devour by \textsc{gulping/bolting} & single piece of food \\
xii.  & -'\'a\'ah & consume \textsc{solid, round object} dunked in liquid & moistened
  bread, cake \\
xiii. & -dl\k{a} & consume/drink \textsc{unspecified} & soup, thin gruel \\
xiv.  & -kaah    & consume/drink from \textsc{open container} & milk from glass \\
xv.   & -t'aah   & consume/drink from \textsc{closed container} & milk from bottle \\
\bottomrule
\end{tabular}
\captionof{table}{Navajo verb stems of consumption (some diacritics simplified).}
\end{center}

\paragraph{Learning languages.} Quotes on the slides: ``To have a second language is to
have a second soul'' (Charlemagne); ``A man who knows four languages is worth four men''
(Charles V); ``I speak English to my accountants, French to my ambassadors, Italian to my
mistress, Latin to my God, and German to my horse'' (Frederick the Great). And the
opposite camp: ``I hate relativism more than I hate anything else, excepting, maybe,
fiberglass powerboats'' (\textbf{Jerry Fodor}, 1935--2017, one of the principal
philosophers of mind of the late twentieth century, very influential on cognitive
science, and a strong critic of linguistic relativity).

\paragraph{Why start from time?}
The most common nouns in English (lists exist also for Italian and French) are, in order:
\textbf{time}, year, people, way, day. Time is the most talked-about thing, so it is a
good place to look for differences between languages.

\subsubsection{Space and time}
\slidesrc{Pragmatics: 31--38}

We cannot see or touch time, so all languages talk about it using \textbf{space}; but
they use space differently.
\begin{itemize}
  \item \textbf{Writing direction}. English speakers arrange a temporal sequence (e.g.\
        pictures of a growing bear, as in the old Nestlé logo) from \emph{left to right};
        speakers of Hebrew or Arabic, who write right to left, arrange it \emph{right to
        left} (Tversky et al., 1991; Fuhrman \& Boroditsky, 2010).
  \item \textbf{Front and back}. For English speakers the future is \emph{ahead} and the
        past \emph{behind}. In \textbf{Aymara} (Andes) the past is \emph{in front} (it
        is known, you can ``see'' it) and the future \emph{behind}: speakers gesture
        forward when talking about the past (Núñez \& Sweetser, 2006).
  \item \textbf{Up and down}. In Mandarin, last week is the ``up'' week and next week the
        ``down'' week: time is also \emph{vertical}. So time can be laid on three axes:
        \textbf{transverse} (left--right), \textbf{sagittal} (front--back),
        \textbf{vertical} (up--down) (Boroditsky et al., 2010; Fuhrman et al., 2011).
  \item \textbf{Absolute directions}. In some Aboriginal Australian communities
        (Kuuk Thaayorre, Pormpuraaw) people do not use ``left'' and ``right'' but
        cardinal directions (``there is an ant on your south-west leg''). Given cards
        showing a man at different ages, they arrange them \textbf{from east to west},
        whichever way they are facing: left to right when facing south, towards their
        body when facing east (Boroditsky \& Gaby, 2010).
  \item \textbf{Uphill and downhill}. The Yupno of Papua New Guinea place the past
        downhill and the future uphill, following the slope of their valley (Núñez et
        al., 2012).
\end{itemize}

\begin{nbfact}{is language the cause?}
How do we know that language \emph{causes} these differences, and does not only
correlate with them (with culture, landscape, \ldots)? \textbf{Teach people new
metaphors for time} (e.g.\ English speakers trained to talk about time vertically) and
see if this changes the way they think. Answer: \textbf{yes}.
\end{nbfact}

\subsubsection{Colour}
\slidesrc{Pragmatics: 39--40}

The colour spectrum is continuous, but languages cut it into categories differently.
English has one basic word, \emph{blue}; Russian has two: \textbf{goluboy} (light blue)
and \textbf{siniy} (dark blue). In an experiment (Winawer, Witthoft, Frank, Wade,
Boroditsky, PNAS) Russian speakers were \textbf{faster at distinguishing two blues when
they fell on different sides} of the goluboy/siniy boundary; English speakers showed no
such advantage. The words we have change, measurably, how we perceive.

\subsubsection{Grammatical gender}
\slidesrc{Pragmatics: 41--44}

In German the sun is feminine (\emph{die Sonne}) and the moon masculine (\emph{der
Mond}); in Spanish it is the opposite (\emph{el sol}, \emph{la luna}). Gender is not a
small detail of a few words: it is \textbf{pervasive}. Russian examples with \emph{okno}
(window, neuter), \emph{taburetka} (stool, feminine) and \emph{stul} (chair,
masculine):

\begin{center}
\small
\begin{tabularx}{\linewidth}{@{}>{\raggedright\arraybackslash}p{3.6cm} X@{}}
\toprule
\textbf{where gender shows up} & \textbf{example} \\
\midrule
different phonology & okn\textbf{o} (-o, -e = neuter), taburetk\textbf{a} (-a =
  feminine), stu\textbf{l} (consonant = masculine) \\
different number words & \textbf{odna} taburetka vs.\ \textbf{odin} stul (one stool,
  one chair) \\
different adjective endings & bel\textbf{oye} okno, bel\textbf{aya} taburetka,
  bel\textbf{y} stul (white window/stool/chair) \\
different pronouns and possessives & \textbf{ono} beloye, \textbf{ona} belaya, \textbf{on}
  bely; \textbf{moyo} okno beloye, \textbf{moya} taburetka belaya, \textbf{moy} stul
  bely \\
different verb endings & moyo okno \textbf{bilo} beloye, moya taburetka \textbf{bila}
  belaya, moy stul \textbf{bil} bely (my \ldots\ was white) \\
\bottomrule
\end{tabularx}
\end{center}
To say ``my chair was white'' a Russian speaker must mark the masculine gender of
\emph{chair} on the possessive, the verb and the adjective.

\textbf{Do people really take these genders as meaningful?} Yes: speakers use
grammatical gender when they give voices to animated characters, when they describe
objects, they underestimate hurricanes with female names, and artists personify abstract
entities according to the gender of the word. Michelangelo's statues on the tomb of
Giuliano de' Medici (1526--1531) and of Lorenzo de' Medici represent \emph{Dawn}
(\emph{l'aurora}) and \emph{Night} (\emph{la notte}) as women, \emph{Day} (\emph{il
giorno}) and \emph{Dusk} (\emph{il crepuscolo}) as men: exactly the genders of the
Italian words.

\subsubsection{Agency: who did it?}
\slidesrc{Pragmatics: 45--49}

\emph{``Pat broke the vase.''} English describes even accidents with an agent. Spanish
distinguishes: \emph{Pat rompió el florero} (Pat broke the vase, on purpose) vs.\
\emph{Se rompió el florero} (the vase broke itself, i.e.\ by accident). Language primes
the same contrast in English: \emph{``The match blew out''} vs.\ \emph{``She blew out the
match''}; \emph{``The toast burned''} vs.\ \emph{``She burned the toast''}.

In the experiment, people watched videos of someone breaking things, intentionally or
by accident, and were later asked \textbf{``who did it the first time?''}. For
intentional events everybody remembers the agent equally well; for \textbf{accidental}
events, English speakers remember the agent better than Spanish (or Japanese) speakers,
whose language would normally describe the event without an agent. What a language makes
you say influences what you notice and remember, and, as studies on agentive vs.\
non-agentive descriptions show, also how much blame you assign.

\begin{nbsummary}[language and thought]
\begin{itemize}[nosep]
  \item Newspeak (Orwell): language designed to make thoughts impossible.
  \item Sapir--Whorf: strong (determinism, false) vs.\ weak (influence, supported by
        evidence).
  \item Languages force different information: tense, gender, aspect, evidentiality;
        Navajo verbs of eating. About 7\,160 languages, 44\% endangered.
  \item Time in space: left--right by writing direction; Aymara past in front; Mandarin
        vertical; Kuuk Thaayorre east-to-west; Yupno uphill future. Teaching new
        metaphors changes thinking $\Rightarrow$ causal.
  \item Colour: Russian goluboy/siniy, faster discrimination across the boundary.
  \item Gender: pervasive (Russian), influences descriptions, voices, art (Michelangelo).
  \item Agency: English agentive descriptions; better memory of agents of accidents.
\end{itemize}
\end{nbsummary}

\subsection{The Language of Numbers and the Power of Names}
\label{sec:numbers-and-names}
\slidesrc{Pragmatics: 50--66}

\subsubsection{A few words on maths}
\slidesrc{Pragmatics: 50--59}

Galileo Galilei (1564--1642), in \emph{Il Saggiatore}, called mathematics a language:

\begin{quote}\small
``Philosophy is written in this grand book, which stands continually open before our
eyes (I say the `Universe'), but cannot be understood without first learning to
comprehend the language and know the characters in which it is written. It is written in
mathematical language, and its characters are triangles, circles and other geometric
figures, without which it is impossible to humanly understand a word; without these one
is wandering in a dark labyrinth.''
\end{quote}

If mathematics is a language, the way numbers are \emph{written and named} should also
matter. The slides follow its history:
\begin{itemize}
  \item \textbf{Greek numerals}. The ancient Greeks first had a system like the Romans;
        from the 4th century BC they used different symbols for 1--9 (as we do), but
        \emph{another} set for 10, 20, 30, \ldots\ and yet another for 100, 200, 300:
        like many ancient systems, they eventually ran out of symbols. They had \textbf{no
        symbol for zero} (\emph{mēdén}) and no digit placeholder: they were unsure
        whether zero was a number at all (``how can nothing be something?''). Yet with
        these numerals \textbf{Eratosthenes} of Cyrene (c.\ 276--194 BC), chief
        librarian of Alexandria, computed the circumference of the Earth with remarkable
        accuracy, invented geography as a discipline, and gave the \emph{sieve of
        Eratosthenes} for finding prime numbers.
  \item \textbf{Roman numerals}. Seven symbols: I $=1$, V $=5$, X $=10$, L $=50$, C
        $=100$, D $=500$, M $=1000$. They remained the usual way of writing numbers in
        Europe well into the Late Middle Ages and were replaced from the 14th century by
        the more convenient Hindu-Arabic numerals; they survive in minor uses today. The
        Romans nevertheless built aqueducts and roads that still stand.
  \item \textbf{Hindu-Arabic numerals}: ten symbols, 1, 2, \ldots, 9, 0, for the
        decimal positional system. They originated in India in the 6th or 7th century
        and reached Europe through Middle Eastern mathematicians (especially
        al-Khwarizmi and al-Kindi) around the 12th century. They were a profound break
        with earlier methods of counting (such as the abacus) and paved the way for
        algebra.
  \item \textbf{Counting with the body}. The Oksapmin of Papua New Guinea count with a
        system of 27 body parts: from the thumb of one hand, up the arm, across the head
        and down the other arm; the number is the name of the body part reached.
  \item \textbf{Languages with few number words}. Some Amazonian languages have no words
        for exact numbers beyond very small ones. In matching tasks (put on the table as
        many objects as you see in a row), speakers succeed with small quantities but
        become approximate with larger ones: without number words, exact large
        quantities are hard to keep in mind.
\end{itemize}

\begin{nbintuition}[the message]
Notation and number words are a language, and some languages make some thoughts easier:
long multiplication is easy in positional decimal notation and painful with Roman
numerals; exact counting is hard without number words. But the Greeks and Romans did
great mathematics and engineering with clumsy numerals. Again: language \emph{influences}
thought (weak Whorf), it does not \emph{determine} it.
\end{nbintuition}

\subsubsection{The importance of what a thing is called}
\slidesrc{Pragmatics: 60--66}

The same thing can be named in ways that push the hearer to judge it differently:
\begin{itemize}[nosep]
  \item \emph{patriots}, \emph{activists}, or \emph{terrorists}?
  \item \emph{torture} or \emph{enhanced interrogation techniques}?
\end{itemize}

\textbf{But sometimes language seems ineffective. Why?} In 2003 the cafeterias of the US
House of Representatives renamed French fries \textbf{``freedom fries''} and French toast
``freedom toast'', a symbolic protest against France's opposition to the war in Iraq
(BBC News, 12 March 2003). It had happened before: during World War~I sauerkraut became
``liberty cabbage'', liverwurst ``liberty sausage'', hamburger ``Salisbury steak'',
dachshunds ``liberty pups''. In 2003 the joke list went on: freedom poodles, freedom
kissing, freedom manicures, and ``France'' itself would become ``freedom land'', ``French''
``the language of freedom''. These renamings did not change much: everybody knew what
the new names stood for, so the old meaning (and the old attitude) remained.

\textbf{And sometimes it works.} Prunes had an image problem (old people, laxatives).
Producers renamed them \textbf{dried plums}, and sales went up: ``Folks favor dried plums
over prunes'', with sales up 5.5\% over the previous year, according to the state
marketing board (Business Journal, 2001). Many other examples from all around the world
concern the names of products, and the classic marketing mistakes when a name has an
unwanted meaning in another language.

\begin{nbwhy}
Naming is pragmatics in action: choosing a word is choosing what the hearer will notice
(an agent or not, an act of violence or a technique, an old person's remedy or a fruit).
NLP systems meet this every day: sentiment and stance detection must recognise that
\emph{terrorist} and \emph{freedom fighter} may refer to the same person with opposite
attitudes, and text generation systems choose words with a pragmatic effect, wanted or
not.
\end{nbwhy}

The deck ends with Nelson Mandela:

\begin{quote}
``If you talk to a man in a language he understands, that goes to his head. If you talk
to him in his language, that goes to his heart.''
\end{quote}

\begin{nbsummary}[pragmatics]
\begin{itemize}[nosep]
  \item Pragmatics: how context (situation, pre-existing knowledge, inferred intent)
        contributes to meaning. Semiotics: the study of signs and meaning-making.
  \item Communication: intentional exchange of information through a shared system of
        signs. Classical view (sentences are true/false) vs.\ modern view (language is
        action).
  \item Wittgenstein (1953): meaning is use in a language-game. Austin (1962): in saying
        something we do something (performatives). Searle (1969): speech acts,
        illocutions, intentionality.
  \item Speech acts: inform, query, command, promise, acknowledge. Stages: intention,
        generation, synthesis (speaker); perception, analysis, disambiguation,
        incorporation (hearer).
  \item Language and thought: Sapir--Whorf strong (false) vs.\ weak (supported):
        time/space, colour, gender, agency, numbers.
  \item Names matter: freedom fries (ineffective), dried plums (effective).
\end{itemize}
\end{nbsummary}

